\chapter{Instruction and Translation}\label{ch:instruction}

Two institutional settings make especially strong use of selective presentation: teaching and translation. In both, a room of people with different needs faces one presentation, and in both, the usual remedies divide the room or slow it to a single pace. This chapter considers what selective streams could do in classrooms and multilingual settings, and where they could do harm.

\section{One demonstration, several explanations}

A shared demonstration, an experiment, a procedure, a worked problem, can be shown to a whole class while each student receives the explanation they need. The depth family of \Cref{ch:depth} is the natural structure: introductory, intermediate, and advanced realizations of one demonstration, synchronized, each complete, with the one-way window allowing any student to move toward more explanation whenever they need it.

The educational value is plain. A single presentation normally forces one pace on everyone, too fast for some and too slow for others. A depth family lets each student follow at their level while remaining in the same room, watching the same demonstration, discussing it afterward with people who saw the same events explained differently. The classroom remains a group.

\section{Temporary scaffolding}

The most useful feature for teaching may be the least conspicuous. A student need not choose a level once. They can remain on the common or advanced realization and step into a more explanatory one for a single difficult interval, returning when they have what they need. The one-way window guarantees that the step toward more explanation is always admissible, and the return becomes admissible once the explanation is complete.

This is scaffolding in the ordinary teaching sense: support taken when it is needed and set aside when it is not, without the student leaving the room, asking in front of others, or being moved to a different group. It is available only because the explanation happens in time the common screen would otherwise have to share, and only because the switch system keeps the student's history coherent when they move.

\section{Selection and tracking}

The same system can be used very differently. An institution could assign students to realizations: this group to the introductory stream, that one to the advanced, on the basis of test scores, prior grades, or a teacher's judgment. That is tracking, whose effects on students assigned to lower tracks have been extensively documented \cite{oakes}, and selective presentation would make it invisible: students in one room, apparently receiving one lesson, actually sorted into levels they did not choose and may not know about.

\Cref{ch:depth} argued that depth should be a property of the presentation a viewer selects, not a rank assigned to the viewer. In a classroom the principle has specific consequences. Students should choose their own realization, with advice if they want it. Their choice should be private: a student's dial position is not the teacher's business unless the student makes it so. They should be able to change it at any time the switch system allows. And nothing that follows should depend on which realization they watched. An assessment that rewards the advanced realization's content penalizes students for choosing help, and turns a tool for learning into a sorting device.

These principles will not always be kept, and the risk should be stated clearly. Selective presentation in schools could make tracking more efficient and less visible than it has ever been. The same apparatus that lets a student take help privately lets an institution withhold material privately. Which it does depends on who controls the dial.

\section{Translation in the room}

\Cref{ch:local} discussed language as a dimension of variation between theaters. In a single multilingual room, the same principles apply at smaller scale. Subtitle streams in several languages cost little light, because they differ only in their text. Explanations can be localized as well as translated, using examples that make sense to each group. Under individual sound, spoken language can vary; under shared sound, the spoken language is common and the others are carried in the picture.

A multilingual classroom or public meeting could thus share one presentation, one room, and one clock while each participant receives it in their own language. The shared score, the shared images, and the rendezvous hold the group together; the language streams let each person take part.

\section{Plural authorship}

Translation at this depth raises a question \Cref{ch:local} began. When a realization is not only translated but localized, with its examples, jokes, and references rewritten, the translator is an author of that realization. In an ordinary film, translation is credited as a service to the original. In a variant family designed with its language realizations from the start, there may be no single original, and each realization's writers are its authors. Credit and provenance should say so: which realization was written by whom, from what source, with what changes. \Cref{ch:versions} develops this into a general requirement.
